% Pursuing Stacks (cotes 134-2 to 134-8): archive cote 134-4, site page 25 (PDF page 26), typescript page 113.
% typescript. Draft compared: ps3.tex:1229-1293.
% AI-transcribed by Claude Opus 5.5 (claude-opus-5-5, adaptive thinking, effort high), Message Batches API, batch ps2 (msgbatch_01NVwr1Yuo8qJf5rykeGBXPK), 2026-11-12; not yet checked by eye.
\scanpage{134-4/25}
\mspage{113}
deduce a notion of \uline{homotopy interval}  i.e.
\[ \hand{\Sigma}_h \hand{\subset} \mathrm{Int}(M) \quad , \]
% NOTE: the right edge of the scan is cut; words at the end of some lines are incomplete (homoto..., ca..., commute...)
namely those intervals $\underline{I}$ in M such that the  two marked sections be homoto%
pic. (NB $\hand{\Sigma}_h$ is non empty if and only if M has a final element,  in this ca%
se $\hand{\Sigma}_h$ contains all intervals such that $\hand{\delta}_o=\hand{\delta}_1$  -  which we may call \uline{trivial}
intervals,  for instance the \uline{final interval} with I=\hand{e}...)  This notion of
a homotopy interval is considerably wider than the one we have worked with
so far,  however it is clearly the right one in the context of \inserted{pure} homotopy
notions.  To avoid any confusion,  we better call this notion by the name
of \uline{weak} \uline{homotopy} \uline{intervals}  -  funnily there won't be any unqualified
\hand{``}homotopy intervals\hand{''} in our present set-up of ``pure'' homotopy notions !

The two prescriptions  above account for twp  among the plain arrows
in our organigram.

\deleted{\illegible{1}}We'll often make use of an \uline{accessory assumption} on R$_h$ ,  which can
be expressed by demanding that \uline{the canonical functor} M$\hand{\to}$ M$_h$  \uline{commute} \uline{to}
\uline{binary products},  in case we assume already such products exist in M.
This can be expressed also by the property that for two maps
\[ f,g \ :\ X \hand{\rightrightarrows} Y_1 x Y_2 \quad , \]
f and g are homotopic iff  so are f$_i$ and g$_i$  ($i\hand{\in\{}1,2\hand{\}}$)  (where the ``only
if'' part is satisfied \xout{anyhow} beforehand anyhow).  This implies \inserted{too} that in
M$_h$ binary products exist\hand{,} \inserted{(and that $f\sim f', g\sim g'$ implies $f\times g\sim f'\times g'$.} \deleted{and th\unsure{at}} On the other hand,  it is trivial that
if M admits a final object,  this is equally a final object of M$_h$  (and
hence,  under the accessory assumption on R$_h$ ,  the functor M$\hand{\to}$M$_h$ commute\unsure{s}
to finite products).

B)  \uline{Homotopisms}.  As a type of structure on M, we got just a subset
\[ W_h \hand{\subset} \mathrm{Fl}(M) \quad , \]
without any basic assumption to make.  The natural \uline{saturation condition}
is just the strong saturation for a subset of Fl(M),  which can be expres-%
sed by stating that W$_h$ can be obtained as the set of arrows made invertib-%
le by some functor from M into a category M',  or equivalently,  by the
localisation functor
\[ \hand{M\to} \deleted{\illegible{1}} W_h^{-1}M \deleted{\illegible{1}} \hand{.} \]
We may refer to a strongly saturated W$_h$ as a ``\uline{homotopism} \uline{structure}''
(or ``homotopy equivalence structure'') in M  -  but as in the case A),
we'll have soon enough to make pretty strong extra assumptions.  Maybe
we should, at the very least,  demand for the notion of homotopy structure
that the canonical functor above,  which is bijective on objects in any
case,  should be moreover \uline{surjective on arrows}  -  thus I'll take this

\margin{25}
